\noindent\textbf{Por qué se usa este procedimiento.} Primero imponer la rigidez $EA/L=k$ y despejar $A=kL/E$; después calcular volumen $AL$ y masa $\rho AL$. La comparación se hace a igual rigidez, no a igual sección.

\begin{enumerate}

\item La masa es $m = \rho k L^2/E$.

\end{enumerate}

\noindent\textbf{Operaciones y resultados intermedios.}

\begin{enumerate}

\item Masa del tirante.

\[\fitmath{y_{1}=\frac{4500\cdot 20000000\cdot {1}^{2}}{1.125\times10^{11}}\approx 0{,}8\,\mathrm{kg}}\]

\end{enumerate}
